\documentclass[a4paper]{article}
\usepackage[a4paper,margin=0mm]{geometry}
\usepackage{fontspec}
\usepackage{unicode-math}
\usepackage{tikz}
\usepackage[table]{xcolor}
\usepackage[absolute]{textpos}
\usepackage{array}

\setmainfont{Libertinus Serif}
\setsansfont{Libertinus Sans}
\setmathfont{Libertinus Math}

\definecolor{accent}{HTML}{0F6A52}
\definecolor{badge}{HTML}{A9CFC2}
\definecolor{ink}{HTML}{2A2E30}

\setlength{\TPHorizModule}{1mm}
\setlength{\TPVertModule}{1mm}
\setlength{\parindent}{0pt}
\pagestyle{empty}


\newcommand{\pagefooter}[1]{%
\begin{textblock*}{194mm}(8mm,286.41mm)
  \begin{tikzpicture}\fill[accent] (0,0) rectangle (194mm,0.6mm);\end{tikzpicture}
\end{textblock*}
\begin{textblock*}{63.59mm}(8mm,288.27mm)
  {\sffamily\bfseries\color{accent}\fontsize{13}{16}\selectfont SAT 1200+}
\end{textblock*}
\begin{textblock*}{10mm}(188mm,288.27mm)
  \centering{\sffamily\color{accent}\fontsize{14}{16}\selectfont #1}
\end{textblock*}}
\newcommand{\frontheader}[1]{%
\begin{textblock*}{194mm}(16mm,16mm)
\begin{tikzpicture}\node[font=\sffamily\bfseries\fontsize{15}{16}\selectfont,text=ink,anchor=west] at (0,0) {#1};\end{tikzpicture}
\end{textblock*}
\begin{textblock*}{194mm}(16mm,22.5mm)
\begin{tikzpicture}\fill[accent] (0,0) rectangle (194mm,0.5mm);\end{tikzpicture}
\end{textblock*}}

\begin{document}
\frontheader{Answer Sheets}
\pagefooter{@@PAGE@@}
\mbox{}
\begin{textblock*}{170mm}(20mm,40mm)
{\fontsize{11}{17}\selectfont\color{ink}
Bu bo'limda 10 ta mock testning javoblari berilgan --- har bir mock uchun bitta sahifa. Faqat ikkala modulni ham tugatgandan keyin oching.

\vspace{4mm}
\textbf{Har bir sahifada:} har bir savol uchun javob, soha (\textbf{ALG} --- Algebra, \textbf{ADV} --- Advanced Math, \textbf{PSDA} --- Problem-Solving \& Data Analysis, \textbf{GT} --- Geometry \& Trigonometry) va daraja (\textbf{E} --- oson, \textbf{M} --- o'rta, \textbf{H} --- qiyin). Yulduzcha (*) bilan belgilangan savollar --- javobni o'zingiz yozadigan (grid-in) savollar.

\vspace{4mm}
\textbf{Xatolar bilan ishlash:}

\vspace{1mm}
\par\vspace{2.2mm}\noindent\hangindent=7mm\hangafter=1 \hspace{1mm}{\color{accent}\textbullet}\hspace{3.5mm}Xato qilgan savollaringizning sohasini sanang: bitta sohada ko'p xato bo'lsa, tegishli darsga qayting.
\par\vspace{2.2mm}\noindent\hangindent=7mm\hangafter=1 \hspace{1mm}{\color{accent}\textbullet}\hspace{3.5mm}1-modulda oson (E) savolda xato qilgan bo'lsangiz --- bu bilim emas, e'tibor muammosi; sekinroq va aniqroq ishlang.
\par\vspace{2.2mm}\noindent\hangindent=7mm\hangafter=1 \hspace{1mm}{\color{accent}\textbullet}\hspace{3.5mm}2-moduldagi qiyin (H) savollarni yecha olmagan bo'lsangiz --- tegishli darsdagi "Summit" to'plamini qayta ishlang.
\par\vspace{2.2mm}\noindent\hangindent=7mm\hangafter=1 \hspace{1mm}{\color{accent}\textbullet}\hspace{3.5mm}Vaqt yetmagan bo'lsa --- bu tezlik muammosi: 5-dars (Speed Tricks) ni takrorlang.\par
}
\end{textblock*}
\end{document}
